\par\Needspace{6\baselineskip}\medskip
\hypertarget{native-143-article104}{}\pdfbookmark[2]{104.}{article-104}
\begin{center}104.\end{center}
Es bleibt, aus $z$, $n$ oder $c$ die Elemente selbst zu bestimmen. Setzt man
\[
\alpha\BareRoot{(ee-1)}=\beta,
\]
\noindent so erhält man aus Gleichung 6 in Art. 99
\[
\mathrm{18)}\qquad\beta=\frac{\sin f\mathpunct{.}\BareRoot{rr'}}{\operatorname{tang}2n}\mathpunct{\cdot}
\]
\noindent Verbindet man diese Formel mit 12 und $12^*$ in Art. 99, so ergibt sich
\[
\begin{aligned}
\mathrm{19)}\qquad&\BareRoot{(ee-1)}=\operatorname{tang}\psi=\frac{\operatorname{tang}f\operatorname{tang}2n}{2(l-z)}\\[.6em]
\mathrm{19^*)}\qquad&\operatorname{tang}\psi=-\frac{\operatorname{tang}f\operatorname{tang}2n}{2(L+z)},
\end{aligned}
\]
\noindent Aus diesen Formeln lässt sich die Exzentrizität bequem und genau berechnen. Die Division von $\beta$ durch $\BareRoot{(ee-1)}$ ergibt $\alpha$, die Multiplikation der beiden Größen ergibt $p$. Damit gilt
\[
\begin{aligned}
\alpha&=\frac{2(l-z)\cos f\mathpunct{.}\BareRoot{rr'}}{\operatorname{tang}2n^2}
=\frac{2mm\cos f\mathpunct{.}\BareRoot{rr'}}{yy\operatorname{tang}2n^2}\\[.7em]
&=\frac{kktt}{4yyrr'\cos f^2\operatorname{tang}2n^2}\\[.7em]
&=\frac{-2(L+z)\cos f\mathpunct{.}\BareRoot{rr'}}{\operatorname{tang}2n^2}
=\frac{-2MM\cos f\mathpunct{.}\BareRoot{rr'}}{YY\operatorname{tang}2n^2}\\[.7em]
&=\frac{kktt}{4YYrr'\cos f^2\operatorname{tang}2n^2}
\end{aligned}
\]
\[
\begin{aligned}
p&=\frac{\sin f\mathpunct{.}\operatorname{tang}f\mathpunct{.}\BareRoot{rr'}}{2(l-z)}
=\frac{yy\sin f\mathpunct{.}\operatorname{tang}f\mathpunct{.}\BareRoot{rr'}}{2mm}\\[.7em]
&=\left(\frac{yrr'\sin2f}{kt}\right)^2\\[.7em]
&=\frac{-\sin f\mathpunct{.}\operatorname{tang}f\mathpunct{.}\BareRoot{rr'}}{2(L+z)}
=\frac{-YY\sin f\mathpunct{.}\operatorname{tang}f\mathpunct{.}\BareRoot{rr'}}{2MM}\\[.7em]
&=\left(\frac{Yrr'\sin2f}{kt}\right)^2\mathpunct{\cdot}
\end{aligned}
\]
\noindent\SourceLeaf{144}Der dritte und der sechste Ausdruck für $p$ stimmen mit den Formeln 18 und $18^*$ in Art. 95 genau überein. Sie zeigen, dass das dort über die Bedeutung der Größen $y$, $Y$ Gesagte auch für die Hyperbel gilt.
\par
Aus der Verbindung der Gleichungen 6 und 9 in Art. 99 folgt
\[
(r'-r)\BareRoot{\frac{ee-1}{rr'}}=e\sin f\mathpunct{.}\left(C-\frac1C\right);
\]
\noindent Führt man also $\psi$ und $\omega$ ein und setzt $C=\operatorname{tang}(45^0+N)$, so erhält man
\[
\mathrm{20)}\qquad\operatorname{tang}2N=\frac{2\sin\psi\operatorname{tang}2\omega}{\sin f\cos2\omega}\mathpunct{\cdot}
\]
\noindent Sobald $C$ auf diese Weise gefunden ist, erhält man für beide Orte die Werte der in Art. 21 mit $u$ bezeichneten Größe. Aus Gleichung III in Art. 21 ergibt sich dann
\[
\begin{aligned}
\operatorname{tang}\tfrac12v&=\frac{C-c}{(C+c)\operatorname{tang}\tfrac12\psi}\\[.6em]
\operatorname{tang}\tfrac12v'&=\frac{Cc-1}{(Cc+1)\operatorname{tang}\tfrac12\psi}
\end{aligned}
\]
\noindent oder wenn man $C$, $c$ durch die Winkel $N$, $n$ ersetzt
\[
\begin{aligned}
\mathrm{21)}\qquad&\operatorname{tang}\tfrac12v=\frac{\sin(N-n)}{\cos(N+n)\operatorname{tang}\tfrac12\psi}\\[.6em]
\mathrm{22)}\qquad&\operatorname{tang}\tfrac12v'=\frac{\sin(N+n)}{\cos(N-n)\operatorname{tang}\tfrac12\psi}\mathpunct{\cdot}
\end{aligned}
\]
\noindent Daraus lassen sich die wahren Anomalien $v$, $v'$ bestimmen. Der Vergleich ihrer Differenz mit $2f$ dient zugleich zur Kontrolle der Rechnung.
\par
Schließlich ergibt sich aus Formel XI in Art. 22 ohne Schwierigkeit für das Zeitintervall vom Perihel bis zum Zeitpunkt, der dem ersten Ort entspricht,
\[
=\frac{\alpha^{\frac32}}k\left\{\frac{2e\cos(N+n)\sin(N-n)}{\cos2N\cos2n}-\log\operatorname{hyp}\frac{\operatorname{tang}(45^0+N)}{\operatorname{tang}(45^0+n)}\right\}
\]
\noindent und ebenso für das Zeitintervall vom Perihel bis zum Zeitpunkt, der dem zweiten Ort entspricht,
\[
=\frac{\alpha^{\frac32}}k\left\{\begin{aligned}&\frac{2e\cos(N-n)\sin(N+n)}{\cos2N\cos2n}\\[.5em]&-\log\operatorname{hyp}\operatorname{tang}(45^0+N)\operatorname{tang}(45^0+n)\end{aligned}\right\}.
\]
\noindent Setzt man also den ersten Zeitpunkt $=T-\tfrac12t$ und demnach den zweiten $=T+\tfrac12t$, so erhält man
\[
\mathrm{23)}\qquad T=\frac{\alpha^{\frac32}}k\left\{\frac{e\operatorname{tang}2N}{\cos2n}-\log\operatorname{hyp}\operatorname{tang}(45^0+N)\right\},
\]
\noindent\SourceLeaf{145}woraus sich der Zeitpunkt des Periheldurchgangs ergibt. Schließlich gilt
\[
\mathrm{24)}\qquad t=\frac{2\alpha^{\frac32}}k\left\{\frac{e\operatorname{tang}2n}{\cos2N}-\log\operatorname{hyp}\operatorname{tang}(45^0+n)\right\},
\]
\noindent Diese Gleichung kann nach Belieben zur abschließenden Kontrolle der Rechnung verwendet werden.
